% Tabla generada automáticamente por 03_tabla_latex.py
% Incluir en la memoria con: % Tabla generada automáticamente por 03_tabla_latex.py
% Incluir en la memoria con: % Tabla generada automáticamente por 03_tabla_latex.py
% Incluir en la memoria con: % Tabla generada automáticamente por 03_tabla_latex.py
% Incluir en la memoria con: \input{tabla_resultados.tex}

\begin{table}[htbp]
\centering
\caption{Resultados experimentales por celda. T\_detect y MTTR expresados como mediana (s). N=10 repeticiones por celda.}
\label{tab:resultados_experimento}
\begin{tabular}{llS[table-format=2.1]S[table-format=3.1]cc}
\toprule
Escenario & Estrategia & {T\_detect med. (s)} & {MTTR med. (s)} & N éxito & Status \\
\midrule
E1 — Cifrado local & A (sin backup) & 0.4 & \textemdash & 0/20 & impossible \\
 & B (Restic+SFTP) & 0.4 & 12.8 & 20/20 & success \\
 & C (MinIO ObjLock) & 0.5 & 12.1 & 20/20 & success \\
\midrule
E2 — Destrucción copias & A (sin backup) & 11.2 & \textemdash & 0/20 & impossible \\
\rowcolor{red!10}
 & B (Restic+SFTP) & 0.7 & \textemdash & 0/20 & impossible \\
\rowcolor{green!15}
 & C (MinIO ObjLock) & 11.1 & 16.8 & 20/20 & success \\
\midrule
E3 — Mov. lateral & A (sin backup) & 1.5 & \textemdash & 0/20 & impossible \\
 & B (Restic+SFTP) & 1.5 & 7.9 & 20/20 & success \\
 & C (MinIO ObjLock) & 1.4 & 12.1 & 20/20 & success \\
\midrule
E4 — Exfiltración & A (sin backup) & 2.0 & \textemdash & 0/20 & impossible \\
 & B (Restic+SFTP) & 1.2 & 10.7 & 20/20 & success \\
 & C (MinIO ObjLock) & 1.2 & 14.8 & 20/20 & success \\
\bottomrule
\end{tabular}
\end{table}

\begin{table}[htbp]
\centering
\caption{Resultados experimentales por celda. T\_detect y MTTR expresados como mediana (s). N=10 repeticiones por celda.}
\label{tab:resultados_experimento}
\begin{tabular}{llS[table-format=2.1]S[table-format=3.1]cc}
\toprule
Escenario & Estrategia & {T\_detect med. (s)} & {MTTR med. (s)} & N éxito & Status \\
\midrule
E1 — Cifrado local & A (sin backup) & 0.4 & \textemdash & 0/20 & impossible \\
 & B (Restic+SFTP) & 0.4 & 12.8 & 20/20 & success \\
 & C (MinIO ObjLock) & 0.5 & 12.1 & 20/20 & success \\
\midrule
E2 — Destrucción copias & A (sin backup) & 11.2 & \textemdash & 0/20 & impossible \\
\rowcolor{red!10}
 & B (Restic+SFTP) & 0.7 & \textemdash & 0/20 & impossible \\
\rowcolor{green!15}
 & C (MinIO ObjLock) & 11.1 & 16.8 & 20/20 & success \\
\midrule
E3 — Mov. lateral & A (sin backup) & 1.5 & \textemdash & 0/20 & impossible \\
 & B (Restic+SFTP) & 1.5 & 7.9 & 20/20 & success \\
 & C (MinIO ObjLock) & 1.4 & 12.1 & 20/20 & success \\
\midrule
E4 — Exfiltración & A (sin backup) & 2.0 & \textemdash & 0/20 & impossible \\
 & B (Restic+SFTP) & 1.2 & 10.7 & 20/20 & success \\
 & C (MinIO ObjLock) & 1.2 & 14.8 & 20/20 & success \\
\bottomrule
\end{tabular}
\end{table}

\begin{table}[htbp]
\centering
\caption{Resultados experimentales por celda. T\_detect y MTTR expresados como mediana (s). N=10 repeticiones por celda.}
\label{tab:resultados_experimento}
\begin{tabular}{llS[table-format=2.1]S[table-format=3.1]cc}
\toprule
Escenario & Estrategia & {T\_detect med. (s)} & {MTTR med. (s)} & N éxito & Status \\
\midrule
E1 — Cifrado local & A (sin backup) & 0.4 & \textemdash & 0/20 & impossible \\
 & B (Restic+SFTP) & 0.4 & 12.8 & 20/20 & success \\
 & C (MinIO ObjLock) & 0.5 & 12.1 & 20/20 & success \\
\midrule
E2 — Destrucción copias & A (sin backup) & 11.2 & \textemdash & 0/20 & impossible \\
\rowcolor{red!10}
 & B (Restic+SFTP) & 0.7 & \textemdash & 0/20 & impossible \\
\rowcolor{green!15}
 & C (MinIO ObjLock) & 11.1 & 16.8 & 20/20 & success \\
\midrule
E3 — Mov. lateral & A (sin backup) & 1.5 & \textemdash & 0/20 & impossible \\
 & B (Restic+SFTP) & 1.5 & 7.9 & 20/20 & success \\
 & C (MinIO ObjLock) & 1.4 & 12.1 & 20/20 & success \\
\midrule
E4 — Exfiltración & A (sin backup) & 2.0 & \textemdash & 0/20 & impossible \\
 & B (Restic+SFTP) & 1.2 & 10.7 & 20/20 & success \\
 & C (MinIO ObjLock) & 1.2 & 14.8 & 20/20 & success \\
\bottomrule
\end{tabular}
\end{table}

\begin{table}[htbp]
\centering
\caption{Resultados experimentales por celda. T\_detect y MTTR expresados como mediana (s). N=10 repeticiones por celda.}
\label{tab:resultados_experimento}
\begin{tabular}{llS[table-format=2.1]S[table-format=3.1]cc}
\toprule
Escenario & Estrategia & {T\_detect med. (s)} & {MTTR med. (s)} & N éxito & Status \\
\midrule
E1 — Cifrado local & A (sin backup) & 0.4 & \textemdash & 0/20 & impossible \\
 & B (Restic+SFTP) & 0.4 & 12.8 & 20/20 & success \\
 & C (MinIO ObjLock) & 0.5 & 12.1 & 20/20 & success \\
\midrule
E2 — Destrucción copias & A (sin backup) & 11.2 & \textemdash & 0/20 & impossible \\
\rowcolor{red!10}
 & B (Restic+SFTP) & 0.7 & \textemdash & 0/20 & impossible \\
\rowcolor{green!15}
 & C (MinIO ObjLock) & 11.1 & 16.8 & 20/20 & success \\
\midrule
E3 — Mov. lateral & A (sin backup) & 1.5 & \textemdash & 0/20 & impossible \\
 & B (Restic+SFTP) & 1.5 & 7.9 & 20/20 & success \\
 & C (MinIO ObjLock) & 1.4 & 12.1 & 20/20 & success \\
\midrule
E4 — Exfiltración & A (sin backup) & 2.0 & \textemdash & 0/20 & impossible \\
 & B (Restic+SFTP) & 1.2 & 10.7 & 20/20 & success \\
 & C (MinIO ObjLock) & 1.2 & 14.8 & 20/20 & success \\
\bottomrule
\end{tabular}
\end{table}